% !TEX program = lualatex
% Prueba 2 de LuaLaTeX en TeXsting: gráficas con datos que calcula Lua, TikZ y MetaPost (luamplib).
\documentclass[11pt,a4paper]{article}
\usepackage[spanish]{babel}
\usepackage{fontspec}
\usepackage{pgfplots}
\pgfplotsset{compat=1.18}
\usetikzlibrary{babel} % para que las flechas -> funcionen con babel en español
\usepackage{luamplib}
\usepackage[margin=2.2cm]{geometry}

% Lua genera los puntos de una respuesta RLC subamortiguada: v(t) = e^{-a t} cos(w t)
\usepackage{luacode}
\begin{luacode*}
  function puntos_rlc(a, w, tmax, n)
    local t = {}
    for i = 0, n do
      local x = tmax * i / n
      t[#t + 1] = string.format("(%.4f,%.4f)", x, math.exp(-a * x) * math.cos(w * x))
    end
    tex.sprint("\\addplot[blue, thick] coordinates {" .. table.concat(t, " ") .. "};")
  end
  function histograma(semilla, n)
    math.randomseed(semilla)
    local c = {}
    for i = 1, 10 do c[i] = 0 end
    for i = 1, n do
      local s = 0
      for k = 1, 6 do s = s + math.random() end
      local b = math.floor(s / 6 * 10) + 1
      if b > 10 then b = 10 end
      c[b] = c[b] + 1
    end
    local t = {}
    for i = 1, 10 do t[#t + 1] = string.format("(%d,%d)", i, c[i]) end
    tex.sprint("\\addplot[fill=teal!60] coordinates {" .. table.concat(t, " ") .. "};")
  end
\end{luacode*}

\title{Prueba 2: gráficas}
\author{TeXsting}
\date{\today}

\begin{document}
\maketitle

\section{Respuesta de un circuito RLC (puntos calculados en Lua)}
\begin{center}
\begin{tikzpicture}
  \begin{axis}[width=0.85\linewidth, height=6cm, xlabel={$t$ (ms)}, ylabel={$v(t)$ (V)},
               grid=major, legend pos=north east]
    \directlua{puntos_rlc(0.5, 6, 6, 150)}
    \addlegendentry{subamortiguado}
    \addplot[red, dashed, domain=0:6, samples=60] {exp(-0.5*x)};
    \addlegendentry{envolvente}
  \end{axis}
\end{tikzpicture}
\end{center}

\section{Histograma (1000 sumas aleatorias en Lua)}
\begin{center}
\begin{tikzpicture}
  \begin{axis}[ybar, width=0.75\linewidth, height=5.5cm, bar width=12pt,
               xlabel={intervalo}, ylabel={frecuencia}, ymin=0]
    \directlua{histograma(42, 1000)}
  \end{axis}
\end{tikzpicture}
\end{center}

\section{Superficie 3D (pgfplots)}
\begin{center}
\begin{tikzpicture}
  \begin{axis}[view={50}{35}, width=9cm, colormap/viridis]
    \addplot3[surf, domain=-2:2, samples=25] {exp(-x^2-y^2)*cos(deg(2*x))};
  \end{axis}
\end{tikzpicture}
\end{center}

\section{TikZ con bucles}
\begin{center}
\begin{tikzpicture}
  \foreach \i in {0,...,11} {
    \pgfmathtruncatemacro{\p}{\i*100/11}
    \draw[fill=red!\p!blue!60] (\i*30:2) circle (0.35);
    \node at (\i*30:2) {\scriptsize \i};
  }
  \draw[thick,->] (0,0) -- (60:1.5);
  \draw[thick,->] (0,0) -- (0:1.2);
\end{tikzpicture}
\end{center}

\section{MetaPost dentro del documento (solo LuaLaTeX)}
\begin{center}
\begin{mplibcode}
beginfig(1);
  u := 1.2cm;
  for i = 0 upto 7:
    draw fullcircle scaled ((i+1)*0.4u) withcolor (i/8)[red, blue];
  endfor;
  path p;
  p := (0,0) for t = 1 upto 120: .. (t/40*u, 0.8u*sind(t*9)) endfor;
  draw p shifted (3.5u, 0) withpen pencircle scaled 1.2pt withcolor 0.2[green, black];
  label.bot("sen(x) en MetaPost", (5u, -1u));
endfig;
\end{mplibcode}
\end{center}
\end{document}
